\documentclass[11pt]{article}
\usepackage[margin=1in]{geometry}
\usepackage{amsmath,amssymb}

\setlength{\parindent}{0pt}
\setlength{\parskip}{0.65\baselineskip}

\newcommand{\Type}[1]{\texttt{#1}}
\newcommand{\Field}[1]{\ensuremath{\mathrm{#1}}}

\title{PuppyCad: A Graph Dialect for Design State}
\author{Draft Specification}
\date{May 9, 2026}

\begin{document}
\maketitle

\begin{abstract}
PuppyCad represents design state as a graph of typed semantic nodes. The graph contains authored intent, constraints, evaluatable features, derived geometry, and diagnostics in one dialect. This document defines the core node form, the initial node vocabulary, the event model, and the derived properties used by evaluators. It does not define a concrete text file format; a textual syntax can lower into the graph dialect described here.
\end{abstract}

\section{Introduction}

Most CAD systems are organized around domain-specific editors and opaque feature histories. Mechanical sketches, solids, PCB traces, schematic nets, and validation messages are often represented in separate internal models connected by ad hoc translation layers. This makes it difficult to preserve design intent, inspect dependencies, or allow software agents to participate in editing at the same semantic level as a human user.

PuppyCad proposes a different foundation. A design is represented as a typed graph. The graph contains authored facts, derived facts, and diagnostic facts. Changes to the graph are recorded as explicit events. Solvers, geometry kernels, validators, and other evaluators do not mutate private internal state as their primary output; instead, they emit graph events that can be inspected, replayed, invalidated, and recomputed.

\section{PCadNode}

All design facts in PuppyCad are \textbf{PCadNodes}. A PCadNode is a record

\[
n = (\Field{id},\ \Field{type},\ \Field{src},\ \Field{arg})
\]

where \Field{id} is stable identity, \Field{type} selects the semantic node kind, \Field{src} contains references to input nodes or subfeatures, and \Field{arg} contains node-local parameters.

\begin{center}
\small
\begin{tabular}{@{}p{0.16\linewidth}p{0.74\linewidth}@{}}
\hline
\textbf{Field} & \textbf{Meaning} \\
\hline
\Field{id} & Stable identifier. The same semantic object should keep the same id across edits. \\
\Field{type} & Node kind, such as \Type{Line}, \Type{Coincident}, \Type{Extrude}, or \Type{Diagnostic}. \\
\Field{src} & Ordered references to source nodes, selectors, or subfeatures. These define graph dependencies. \\
\Field{arg} & Immediate data for the node: dimensions, transforms, flags, names, operation kinds, tolerances. \\
\hline
\end{tabular}
\end{center}

A PuppyCad graph is

\[
G = (N,\ R,\ L)
\]

where $N$ is a set of PCadNodes, $R$ is a set of root node ids, and $L$ is an event log whose rewrites produce the current folded graph state.

Every node type listed in the following subsections is a specialization of PCadNode. The subsection names group node types by role or domain; they do not introduce separate object models.

The node vocabulary is open. Implementations may add domain-specific node types, but new types must preserve stable identity, explicit sources, and inspectable arguments.

\subsection{Terminology}

The following terms are used throughout the node tables.

\begin{center}
\small
\begin{tabular}{@{}p{0.22\linewidth}p{0.68\linewidth}@{}}
\hline
\textbf{Term} & \textbf{Meaning} \\
\hline
Authored & Entered by a user, agent, importer, or external system as design intent. \\
Derived & Produced by an evaluator from source nodes. \\
Source & A node, selector, or subfeature read by another node or evaluator. \\
Input & A direct source listed in \Field{src}. \\
Selector & A reference to part of a node, such as an endpoint, edge, face, pad, or connector frame. \\
Authority & The producer responsible for an event or derived node. \\
Stable identity & An id that remains attached to the same semantic object across edits. \\
Stale & A derived node whose recorded sources no longer match the current graph. \\
Diagnostic & A graph node describing a problem, warning, status, or validation result. \\
Root & A node selected as visible, exported, or evaluated output of the graph. \\
\hline
\end{tabular}
\end{center}

\subsection{Reference Nodes}

References connect PCadNodes without requiring every graph edge to be its own node. A reference resolves against a graph state:

\[
\operatorname{resolve}_G(r) \in \{\text{valid},\ \text{missing},\ \text{stale},\ \text{ambiguous}\}.
\]

\begin{center}
\small
\begin{tabular}{@{}p{0.16\linewidth}p{0.20\linewidth}p{0.22\linewidth}p{0.32\linewidth}@{}}
\hline
\textbf{Type} & \textbf{src} & \textbf{arg} & \textbf{Semantics} \\
\hline
\Type{Ref} & () & target id, role & Reference to a whole node. \\
\Type{Selector} & ref & path, role & Reference to a subfeature, such as endpoint, edge, face, pad, or frame. \\
\Type{NamedRef} & ref & name & Stable human or agent-facing reference name. \\
\Type{QueryRef} & root & predicate & Late-bound reference selected by a query over graph facts. \\
\Type{ExternalRef} & () & uri, kind & Reference to external data, imported geometry, supplier data, or documents. \\
\hline
\end{tabular}
\end{center}

Selectors should be semantic, not positional. Names such as \Type{outerLoop} or \Type{holeAxis} are preferable to raw list indexes into generated topology.

\subsection{Source Nodes}

Source nodes introduce durable design scope, coordinates, quantities, and authored values.

\begin{center}
\small
\begin{tabular}{@{}p{0.18\linewidth}p{0.20\linewidth}p{0.22\linewidth}p{0.30\linewidth}@{}}
\hline
\textbf{Type} & \textbf{src} & \textbf{arg} & \textbf{Semantics} \\
\hline
\Type{Document} & () & name & Top-level design document. \\
\Type{Part} & () & name & Reusable mechanical design scope. \\
\Type{Assembly} & parts? & name & Scope containing instances and mates. \\
\Type{Board} & () & name, stackup? & PCB or electronics design scope. \\
\Type{Param} & expr? & name, unit, value? & Named design value. May be literal or expression-derived. \\
\Type{Const} & () & value, unit? & Literal scalar, vector, string, boolean, or enum. \\
\Type{Quantity} & value & unit & Numeric value with dimensional unit. \\
\Type{Expr} & inputs & expression & Derived scalar or vector expression over parameters. \\
\Type{Frame} & parent? & transform & Coordinate frame in 2D or 3D space. \\
\Type{Material} & () & properties & Material assignment or material definition. \\
\hline
\end{tabular}
\end{center}

Quantities carry units. Evaluators must normalize compatible units before comparison or computation.

\subsection{Sketch Nodes}

Sketch nodes describe two-dimensional authored geometry in a local frame. Sketch geometry is intent; profiles and solved placements are derived facts.

\begin{center}
\small
\begin{tabular}{@{}p{0.18\linewidth}p{0.20\linewidth}p{0.22\linewidth}p{0.30\linewidth}@{}}
\hline
\textbf{Type} & \textbf{src} & \textbf{arg} & \textbf{Semantics} \\
\hline
\Type{Sketch} & frame & name, policy & 2D modeling context on a frame or workplane. \\
\Type{Point} & sketch & x, y & Sketch point in local coordinates. \\
\Type{Line} & point, point & construction? & Segment between two points. \\
\Type{Circle} & center & radius & Circle carrier in sketch coordinates. \\
\Type{Arc} & center, start, end & direction? & Circular arc. \\
\Type{Spline} & points & degree, knots? & Freeform sketch curve. \\
\Type{Text} & sketch & content, pose, style & Text entity that may become outline geometry. \\
\Type{Projection} & source & target sketch, selector & Projected reference geometry. \\
\Type{Profile} & loops & policy & Closed region selected or detected for features. \\
\hline
\end{tabular}
\end{center}

Construction geometry participates in constraints and references but is excluded from material-producing profiles unless explicitly selected.

\subsection{Constraint Nodes}

Constraint nodes are authored relations over references. A constraint declares what should hold; an evaluator decides whether it holds and may emit repair rewrites or diagnostics.

\[
c = (\Field{kind},\ \Field{src},\ \Field{arg},\ \Field{policy})
\]

\begin{center}
\small
\begin{tabular}{@{}p{0.19\linewidth}p{0.22\linewidth}p{0.20\linewidth}p{0.29\linewidth}@{}}
\hline
\textbf{Type} & \textbf{src} & \textbf{arg} & \textbf{Semantics} \\
\hline
\Type{Coincident} & ref, ref & tolerance? & Resolved points or carriers occupy the same location. \\
\Type{Horizontal} & line & tolerance? & Line direction is horizontal in sketch frame. \\
\Type{Vertical} & line & tolerance? & Line direction is vertical in sketch frame. \\
\Type{Parallel} & line, line & tolerance? & Lines have equal direction up to sign. \\
\Type{Perpendicular} & line, line & tolerance? & Lines meet at a right angle. \\
\Type{Tangent} & curve, curve & tolerance? & Curves share tangent direction at contact. \\
\Type{Concentric} & circle, circle & tolerance? & Circular carriers share a center. \\
\Type{Equal} & ref, ref & measure & Referenced measures are equal. \\
\Type{Distance} & ref, ref & quantity, mode & Separation equals or satisfies the target quantity. \\
\Type{Angle} & ref, ref & quantity & Relative angle equals target. \\
\Type{Radius} & circle/arc & quantity & Radius equals target. \\
\Type{Diameter} & circle/arc & quantity & Diameter equals target. \\
\Type{Fixed} & ref & dof? & Selected degrees of freedom are not editable. \\
\Type{Symmetric} & ref, ref, axis & tolerance? & Two entities are symmetric across axis. \\
\hline
\end{tabular}
\end{center}

Constraint status is one of \Type{satisfied}, \Type{unsatisfied}, \Type{undefined}, \Type{ambiguous}, \Type{overconstrained}, or \Type{stale}. Constraint policy records whether the relation is \Type{driving}, \Type{checked}, or \Type{advisory}.

\subsection{Feature Nodes}

Feature nodes transform references and parameters into derived geometry. Feature order may be represented by data dependencies, explicit ordering references, or both.

\begin{center}
\small
\begin{tabular}{@{}p{0.18\linewidth}p{0.22\linewidth}p{0.22\linewidth}p{0.28\linewidth}@{}}
\hline
\textbf{Type} & \textbf{src} & \textbf{arg} & \textbf{Semantics} \\
\hline
\Type{Extrude} & profile & distance, direction, operation & Sweep profile linearly to create, join, or cut body. \\
\Type{Revolve} & profile, axis & angle, operation & Rotate profile around axis. \\
\Type{Sweep} & profile, path & orientation policy & Sweep profile along path. \\
\Type{Loft} & profiles & continuity policy & Interpolate body through profiles. \\
\Type{Fillet} & edge refs & radius, continuity & Round selected edges. \\
\Type{Chamfer} & edge refs & distance or angle & Bevel selected edges. \\
\Type{Shell} & body, face refs & thickness & Hollow body by offsetting faces. \\
\Type{Pattern} & source & direction/count or transform set & Repeat feature or body. \\
\Type{Mirror} & source, plane & keep original? & Reflect feature or body across plane. \\
\Type{Boolean} & bodies & union, subtract, intersect & Combine solid bodies. \\
\Type{Suppress} & feature & reason? & Keep intent while disabling normal outputs. \\
\hline
\end{tabular}
\end{center}

Feature evaluation emits derived topology and diagnostics. Failed or stale features should preserve authored intent and report their failure as graph facts.

\subsection{Topology Nodes}

Topology nodes describe generated or imported shape structure. They may be authored in imported geometry, but they are usually evaluator outputs.

\begin{center}
\small
\begin{tabular}{@{}p{0.18\linewidth}p{0.22\linewidth}p{0.22\linewidth}p{0.28\linewidth}@{}}
\hline
\textbf{Type} & \textbf{src} & \textbf{arg} & \textbf{Semantics} \\
\hline
\Type{Body} & features? & kind, validity & Solid, surface, mesh, or wire body. \\
\Type{Shell} & faces & closed?, orientation & Connected set of oriented faces. \\
\Type{Face} & surface, loops & orientation & Bounded surface region. \\
\Type{Loop} & edges & outer? & Ordered boundary loop. \\
\Type{Edge} & curve, vertices & orientation & Bounded curve in topology. \\
\Type{Vertex} & point & tolerance? & Topological point. \\
\Type{Curve} & control data & kind & Geometric curve carrier. \\
\Type{Surface} & control data & kind & Geometric surface carrier. \\
\Type{Mesh} & body? & vertices, faces & Tessellated representation for display or analysis. \\
\hline
\end{tabular}
\end{center}

Generated topology must carry provenance sufficient for downstream selectors. A downstream reference should target semantic evidence, not creation order.

\subsection{Assembly Nodes}

Assembly nodes place reusable definitions into spatial relationships.

\begin{center}
\small
\begin{tabular}{@{}p{0.18\linewidth}p{0.22\linewidth}p{0.22\linewidth}p{0.28\linewidth}@{}}
\hline
\textbf{Type} & \textbf{src} & \textbf{arg} & \textbf{Semantics} \\
\hline
\Type{Instance} & part or assembly & transform, overrides & Occurrence of reusable design content. \\
\Type{Connector} & instance or part & local frame, role & Named attachment frame. \\
\Type{Mate} & connector, connector & kind, limits? & Spatial constraint between connectors. \\
\Type{Ground} & instance & transform? & Removes global rigid motion for an instance. \\
\Type{JointLimit} & mate & min, max & Motion bounds for a joint. \\
\Type{ExplodeStep} & instances & transform offsets & Presentation state for exploded views. \\
\hline
\end{tabular}
\end{center}

Assembly solvers produce resolved poses, degrees of freedom, and diagnostics as derived nodes.

\subsubsection{Mate Kinds}

Mate kinds constrain two connector frames. The vocabulary is open, but the following kinds define the initial mechanical connector set.

\begin{center}
\small
\begin{tabular}{@{}p{0.18\linewidth}p{0.22\linewidth}p{0.22\linewidth}p{0.28\linewidth}@{}}
\hline
\textbf{Kind} & \textbf{src} & \textbf{arg} & \textbf{Semantics} \\
\hline
\Type{fasten} & connector, connector & offset? & Connector frames coincide; no relative degrees of freedom remain. \\
\Type{revolute} & connector, connector & axis, limits? & Origins coincide and axes align; rotation about the axis remains free. \\
\Type{prismatic} & connector, connector & axis, limits? & Axes align; translation along the axis remains free. \\
\Type{cylindrical} & connector, connector & axis, limits? & Axes align; rotation and translation along the axis remain free. \\
\Type{planar} & connector, connector & normal, limits? & One frame remains on a plane; two in-plane translations and one normal rotation remain free. \\
\Type{ball} & connector, connector & limits? & Origins coincide; orientation remains free. \\
\Type{parallel} & connector, connector & axis refs & Selected axes remain parallel. \\
\Type{perpendicular} & connector, connector & axis refs & Selected axes remain perpendicular. \\
\Type{tangent} & connector or face refs & side, offset? & Referenced surfaces or curves remain tangent. \\
\Type{distance} & connector or face refs & quantity & Separation equals or satisfies the target distance. \\
\Type{angle} & connector or axis refs & quantity & Relative angle equals or satisfies the target angle. \\
\Type{gear} & revolute, revolute & ratio, phase? & Rotations are coupled by a gear ratio. \\
\Type{rackPinion} & revolute, prismatic & ratio, phase? & Rotation is coupled to linear travel. \\
\hline
\end{tabular}
\end{center}

\subsection{Electronics Nodes}

Electronics and PCB objects use the same graph form. These nodes are optional domain vocabulary, not a separate model.

\begin{center}
\small
\begin{tabular}{@{}p{0.18\linewidth}p{0.22\linewidth}p{0.22\linewidth}p{0.28\linewidth}@{}}
\hline
\textbf{Type} & \textbf{src} & \textbf{arg} & \textbf{Semantics} \\
\hline
\Type{Component} & footprint?, symbol? & value, supplier? & Logical and physical component identity. \\
\Type{Footprint} & pads & package data & Physical land pattern. \\
\Type{Pad} & footprint & shape, layer, net? & Conductive attachment region. \\
\Type{Net} & pins or pads & name, class & Electrical connectivity group. \\
\Type{Trace} & path & width, layer, net & Routed copper path. \\
\Type{Via} & point & drill, layers, net & Layer transition. \\
\Type{Zone} & boundary & layer, net, clearance & Filled copper region. \\
\Type{Keepout} & region & layers, rule & Prohibited placement or routing area. \\
\hline
\end{tabular}
\end{center}

Mechanical and electrical domains may share frames, constraints, and references. A board outline, connector frame, and enclosure feature can therefore participate in one graph.

\subsection{Evaluator Nodes}

Evaluator nodes are derived facts produced by solvers, kernels, validators, and analyzers. They are graph facts, not hidden caches.

\begin{center}
\small
\begin{tabular}{@{}p{0.20\linewidth}p{0.22\linewidth}p{0.22\linewidth}p{0.26\linewidth}@{}}
\hline
\textbf{Type} & \textbf{src} & \textbf{arg} & \textbf{Semantics} \\
\hline
\Type{SolvedSketch} & sketch, constraints & status, residuals & Solver result for a sketch. \\
\Type{ProfileResult} & sketch entities & loops, regions & Detected closed profiles. \\
\Type{GeneratedBody} & feature & body refs & Body produced by feature evaluation. \\
\Type{ResolvedPose} & instance, mates & transform & Assembly instance pose. \\
\Type{AssemblyDOF} & instances, mates & count, basis & Derived assembly degrees of freedom. \\
\Type{MassProperties} & body, material? & volume, area, center, inertia & Physical properties. \\
\Type{CollisionResult} & bodies or instances & contacts & Collision or interference result. \\
\Type{SimulationResult} & model, loads & fields, convergence & Analysis output. \\
\Type{Diagnostic} & sources & severity, code, message & Validation or evaluation observation. \\
\hline
\end{tabular}
\end{center}

Evaluator events should record \Type{authority}, \Type{authorityVersion}, \Type{sourceIds}, and \Type{sourceHash}. If a source changes, the output is stale until recomputed or explicitly accepted.

\section{Events and Rewrites}

Events are the durable mutation mechanism. The folded graph state is produced by applying event rewrites in order.

\[
G_0 + e_1 + e_2 + \cdots + e_n \rightarrow G_n
\]

\begin{center}
\small
\begin{tabular}{@{}p{0.20\linewidth}p{0.24\linewidth}p{0.46\linewidth}@{}}
\hline
\textbf{Rewrite} & \textbf{arg} & \textbf{Semantics} \\
\hline
\Type{AddNode} & node & Insert a new PCadNode. The id must not already be live. \\
\Type{ReplaceNode} & id, node & Replace a node revision while preserving stable identity. \\
\Type{RemoveNodes} & ids & Remove live nodes. Dependent nodes become missing or stale. \\
\Type{SetRootNodes} & ids & Define visible, exported, or evaluated graph roots. \\
\Type{MarkStale} & ids, reason & Preserve nodes but mark them invalid relative to sources. \\
\hline
\end{tabular}
\end{center}

\begin{center}
\small
\begin{tabular}{@{}p{0.20\linewidth}p{0.70\linewidth}@{}}
\hline
\textbf{Event Field} & \textbf{Meaning} \\
\hline
\Field{id} & Stable event id. \\
\Field{authority} & Producer such as user, agent, solver, kernel, router, importer, validator. \\
\Field{rewrites} & Ordered list of graph rewrites. \\
\Field{sourceIds} & Nodes consulted while producing the event. \\
\Field{sourceHash} & Digest of source content relevant to staleness. \\
\Field{timestamp} & Optional wall-clock record, not part of semantic ordering. \\
\hline
\end{tabular}
\end{center}

\section{Derived Properties}

Derived properties are computed from node fields and event history. They are not independent authoring channels.

\begin{center}
\small
\begin{tabular}{@{}p{0.20\linewidth}p{0.70\linewidth}@{}}
\hline
\textbf{Property} & \textbf{Rule} \\
\hline
\Field{domain} & From node type: core, sketch, solid, assembly, pcb, drawing, simulation. \\
\Field{unit} & From quantities, parameters, and expression operands. Compatible units normalize before evaluation. \\
\Field{frame} & Inherited from source sketch, body, instance, or explicit \Type{Frame}. \\
\Field{dependencies} & Transitive closure of \Field{src} references and evaluator-recorded source ids. \\
\Field{validity} & Valid when references resolve and source hashes match; otherwise missing, stale, ambiguous, or invalid. \\
\Field{authority} & Producer class from the event record. Authored and derived facts are distinguished by authority. \\
\Field{diagnostics} & Diagnostics whose sources include this node. \\
\hline
\end{tabular}
\end{center}

\section{Validity}

A PuppyCad graph is valid for a chosen evaluation target when:

\begin{enumerate}
  \item every root id references a live node,
  \item required references resolve to valid nodes or valid selectors,
  \item quantities have compatible units at each operation,
  \item constraints are evaluated under an explicit policy,
  \item derived nodes either match their recorded source hashes or are marked stale, and
  \item diagnostics are represented as PCadNodes rather than hidden evaluator state.
\end{enumerate}

Validity does not require every constraint to be satisfied or every feature to succeed. Unsatisfied constraints, failed features, invalid topology, and stale evaluator outputs are valid graph states when represented by diagnostics.

\section{Execution}

The minimal execution cycle is:

\begin{enumerate}
  \item apply authored or imported events,
  \item resolve references from changed nodes,
  \item run evaluators over affected subgraphs,
  \item emit derived nodes and diagnostics as events,
  \item mark or remove stale outputs when sources change.
\end{enumerate}

This cycle keeps design intent, computed consequences, and validation output inspectable through the same graph dialect.

\end{document}
